% !PRINTED 273
\weberpage{273}

%<<CONT
findet vor und nach dem Durchgang von $f_\nu$ durch Null immer
ein Zeichenwechsel statt, weil $f_{\nu+1}$ und $f_{\nu-1}$
entgegengesetzte Zeichen haben.

Wenn aber $f(x)$ durch Null geht, so geht beim Durchgang durch
$x$ wegen 4. ein Zeichenwechsel verloren. Denn geht $f(x)$ von
negativen zu positiven Werthen, so ist (vgl. \S.\,32) $f'(x)$
und mithin $f_1(x)$ positiv und die Vorzeichen ändern sich so:
\[
\begin{array}{c}
f(x), \; f_1(x) \\
- \quad + \\
+ \quad +,
\end{array}
\]
und wenn $f(x)$ von positiven zu negativen Werthen übergeht, so
ist $f'(x)$ und $f_1(x)$ negativ, also die Zeichen so:
\[
\begin{array}{c}
f(x), \; f_1(x) \\
+ \quad - \\
- \quad -,
\end{array}
\]
mithin ist in beiden Fällen ein Zeichenwechsel verloren gegangen.

Da nun $f_m(x)$ nach 2. sein Zeichen nicht wechselt, so folgt,
dass $V(\alpha) - V(\beta)$ gleich der Anzahl der zwischen
$\alpha$ und $\beta$ gelegenen Wurzeln von $f(x) = 0$ ist, oder:

\gesperrt{Die Anzahl der Wurzeln von $f(x) = 0$ zwischen $\alpha$ und
$\beta$ ist gleich dem Ueberschuss der Anzahl der
Zeichenwechsel der Kette für $x = \alpha$ über die Anzahl der
Zeichenwechsel für $x = \beta$.}

Wenn für $x = \alpha$ oder $x = \beta$ eine oder einige der
mittleren Functionen von $(\mathfrak{R})$ verschwinden sollten,
so ist es wegen 3. gleichgültig, ob wir diesen verschwindenden
Werth durch einen positiven oder einen negativen ersetzen.

\begin{center}
\S.\,87.

\gesperrt{Erstes Beispiel: Kugelfunctionen.}
\end{center}

Ehe wir nun zu den allgemeinen Methoden übergehen, nach denen
Sturm'sche Ketten zu bilden sind, wollen wir einige Beispiele
betrachten, in denen sich durch besondere Umstände solche Ketten
darbieten.
